\chapter{Of the Department of Mines}
\label{art:mines}

\articlenote{In which is set down the charge of the Department of Mines: the winning of ore from the depths below the silo, the limits set upon the miners' working, and the oversight of the silo's ore supply for the working of metal and the silo's own repair.}

\section{The Charge of Mines}

The Department of Mines is charged with the winning of ore from the workings below the Down Deep, that the silo may have metal for its own working and repair, and that the silo's reserves of ore may never fall below such quantity as the Head of Mines judges necessary for the silo's keeping.

\begin{clauses}
\item The Department is led by a Head of Mines, raised up from among its own trades by the shadowing entered into under Article 13, and confirmed in office by the Mayor.
\item Every citizen assigned to a trade of Mines serves under the Head of Mines in the ordinary conduct of that trade, though such a citizen remains bound in every other respect to this Pact, to the Sheriff, and to Judicial exactly as any other citizen is bound.
\item The Head of Mines shall reside and keep the Department's offices in the Down Deep, among the workings described in this Article, and not Up Top.
\end{clauses}

\section{The Winning of Ore}

The removal and processing of ore from the silo's depths is a labor of great weight, requiring strength of body, steadiness of nerve, and intimate knowledge of the workings' structure and history. This section sets down the scope and method of the Department's charge.

\begin{clauses}
\item The winning of ore in the workings below the Down Deep is a standing charge of Mines and needs no order under Article 9 for its ordinary continuation.
\item Every citizen assigned to the winning of ore works under the direct supervision of a miner of proven skill, and no citizen may work alone in any working, regardless of experience or office.
\item The ore won shall be brought to the processing stations under the charge of Mines, where it is measured, recorded, and delivered to the Department of Mechanical for the working of metal and the silo's repair under Article 9.
\item No citizen not assigned to Mines may enter any working or remove any ore or stone from below the Down Deep save under direct written order of the Head of Mines, countersigned by the Mayor, for a purpose set down in that order.
\item A citizen sentenced to a term in the mines under Article 17 labors in the workings under the direction of the Head of Mines for the span the Judge has set, and is counted for that span as assigned to Mines for every purpose of this Article, though not of Article 13; such a citizen works under the same supervision, the same shoring, and the same limit under Section 3 as any miner, and is lodged in the Down Deep under Article 14 for the term's duration. The Head of Mines shall report to Judicial, at each turn of the season, the labor and condition of every citizen so sentenced, and the day each term ends.
\end{clauses}

\section{The Two-Hundred-Foot Radius}

The founders, in their engineering of the silo, set a limit upon how far the miners might work outward from the silo's own structure. This limit stands as their judgment and is not to be questioned or exceeded.

\begin{clauses}
\item No working of ore may run, in any direction, beyond two hundred feet from the silo's own structure, for cause the builders judged sufficient and did not set down.
\item Any citizen who works ore, or directs the working of ore, beyond this limit is treated as among the gravest offenses under Article 17, without regard to ignorance of the limit or the purpose of the work.
\item The Head of Mines is charged with knowing the silo's structure and the distance it is, and with ensuring that no working extends beyond the limit. A citizen assigned to Mines who believes a working has been carried, or is being carried, beyond the limit may report the matter directly to Judicial, without first carrying it to the Head of Mines, and Judicial shall inquire of the Head of Mines and rule upon the working's lawfulness.
\end{clauses}

\section{Coordination with Mechanical and Judicial}

The Department of Mines works in the Down Deep alongside the Department of Mechanical, and both must keep the other informed of the work being done, that neither undermines the silo's structure or the other's safety.

\begin{clauses}
\item The Head of Mines and the Head of Mechanical shall meet no fewer than once per ten days to discuss the workings being opened, the shoring being done, and the state of the structures shared between Mechanical and Mines.
\item Where a working threatens the integrity of the silo's structure or the water lines, the generator's housing, or any other charge of Mechanical, the Head of Mechanical may order the working halted until the threat is resolved, and may carry that order to Judicial if the Head of Mines refuses to comply.
\item Where a citizen assigned to Mines is suspected of taking ore beyond the two-hundred-foot limit, or of deliberate failure to measure the distance, Judicial shall investigate and make a finding under Article 6 before any punishment is carried out.
\end{clauses}

\section{The Trades of Mines}

Every trade of Mines --- the breaking and removal of ore, the processing and measuring of ore, the shoring and tending of workings, and the keeping of records and accounting --- is entered as every other trade is entered, by the shadowing set down under Article 13, save that no citizen may shadow into the breaking of ore or the measuring of distance without also being found fit by the Head of Mines, beyond whatever fitness Article 13 alone would require.

\begin{clauses}
\item The Head of Mines shall keep, and pass to their own shadow before leaving office, a full accounting of the ore won by season, the workings that have been opened and closed, and the distances measured from the silo's structure to the farthest working.
\item No citizen shall be assigned to Mines against the citizen's own naming of a trade at the age Article 13 sets down, save where the Department's need for a given trade cannot otherwise be met, in which case Judicial shall confirm that the assignment is not made in punishment.
\end{clauses}

\section{The Shifts and the Shoring}

The workings are kept by shifts, and the shoring is kept before the ore.

\begin{clauses}
\item The workings are worked in two shifts in every day, and no citizen, whether of the trade or sentenced under Article 17, works below for more than one shift in a day, nor for more than the day's shift on nine days of the cycle, the tenth being the rest day under Article 1; the Head of Mines posts the shifts at the Department's door at each turn of the season.
\item No working is broken or advanced until the shoring of it is set and the shorer of the shift has entered in the shift's log that it is sound; a citizen who breaks ore beyond the last set shoring is answerable under Article 17, and the miner of proven skill under Section 2 who allowed it with them.
\item At the change of every shift, the miners going off and coming on count each other at the head of the workings, and the count is entered in the shift's log with the names of every citizen below; no shift goes off until every citizen of it is counted up, and a citizen missing at the count is sought before any other thing is done.
\item The shorers of Mines are a trade of full shadowing under Section 5, and a shorer's judgment that a working is unsound stands against any order to advance it, from the Head of Mines or any officer, until the Head of Mechanical under Section 4 has seen the working and entered otherwise.
\end{clauses}

\section{The Measuring of Distance}

The limit of Section 3 is kept by measuring, and the measuring is kept by more hands than one.

\begin{clauses}
\item The Head of Mines keeps a measurer for every working, a citizen of the trade found fit under Section 5, who measures the working's farthest face from the silo's own structure at the start of every cycle, by a chain of a length the Head of Mechanical has proved and marked, and enters the distance in the working's log with the day and the measurer's name.
\item The chain of measure is kept at the head of the workings, and no other is used; it is proved against the marks of the silo's own structure once in every season by the Head of Mines and the Head of Mechanical together, and the proving entered.
\item Where a working's measured distance reaches nine parts in ten of the limit under Section 3, the measurer enters it so, the Head of Mines is told the same day, and the working is measured at the start of every shift thereafter until it is closed or turned; a working that reaches the limit is closed and its face shored blind, and the closing entered.
\item The measurer of a working and the miner of proven skill under Section 2 are never the same citizen; and any citizen of Mines may ask to see a working's log of distance and may not be refused.
\end{clauses}

\section{Accident in the Workings}

The workings fall, and flood, and choke; a Department that digs for the silo keeps the reckoning of what the digging costs.

\begin{clauses}
\item A collapse, a flooding, or a fouling of the air in any working is met before all else by the counting-up of every citizen below under Section 6; the Head of Mines is told at once, and the Head of Mechanical within the same shift, and the two order together what the shoring, the pumps, and the vents require.
\item A citizen hurt in the workings is carried to the Down Deep's aid station under Article 19 before any other thing is done; a death is reported to the Infirmary and to Judicial the same day, and the place of it is left as it lies until an Officer of Judicial has seen it, so far as the safety of those still below allows.
\item The Head of Mines keeps a roll of accidents as Article 9 provides for Mechanical, reads it to the Department at each turn of the season, copies it to the Judicial archive, and reports its count yearly to the Mayor; where a collapse came of a working advanced beyond its shoring, or beyond the limit under Section 3, the Head reports it to Judicial as an offense.
\item A citizen sentenced under Article 17 who is hurt in the workings is cared for by the Infirmary exactly as a citizen of the trade, and the term does not run while the citizen cannot work below; a citizen so hurt who cannot return below serves what remains of the term in the labor of the processing stations under Section 2, or, where the Head Physician finds the citizen unfit for that, petitions Judicial to shorten the term, and Judicial hears the petition as Article 6 provides.
\end{clauses}
